\section*{Appendix}
\addcontentsline{toc}{section}{Appendix}
\label{sec:appendix}

% Optional supporting material ONLY. Remember: the brief counts ALL pages
% toward the 15-page limit, so use the appendix sparingly.

% Examples you may include if space allows:
% - Full distribution-fit screenshots (histogram + overlay + KS).
% - Full factorial results table.
% - AnyLogic block-level screenshots.

\subsection*{Software artefacts}

\textbf{AnyLogic model.} \texttt{FermaCore\_Replication\_Harness\_final.alp}
is the model used for every result. It implements the full process flow
(Sec.~\ref{sec:conceptual}), the press scrap, a selectable dryer-queue rule
(\texttt{queueStrategy}) and a configurable order-book size
(\texttt{targetBatchCount}), and a Parameter-Variation experiment
(\texttt{ReplicationsExp}) that runs $n=20$ replications and writes one KPI row
per replication to CSV. A scenario is run by setting the \texttt{Main} parameters
(e.g.\ \texttt{numDryers}, \texttt{cleaningTeamCapacity}, \texttt{targetBatchCount},
\texttt{pressScrapRate}) and re-running the experiment.

\textbf{Analysis scripts (Python).}
\begin{itemize}\setlength\itemsep{0pt}
  \item \texttt{halfwidth\_analysis.py}: replication sizing (half-width method, required $n$ per KPI) for one scenario CSV.
  \item \texttt{scenario\_comparison.py}: cross-scenario KPI table with \SI{95}{\percent} CIs, two-sample significance versus baseline, and the $2^3$ factorial effects.
  \item \texttt{make\_figures.py}: builds the throughput, utilisation, DOE-Pareto and volume figures from \texttt{data/}.
  \item \texttt{make\_fit\_figures.py}: builds the three input-distribution fit figures from \texttt{FermaCore\_input\_data.xlsx}.
\end{itemize}
